\section{Memoria a largo plazo (Long-term Memory)}

\subsection{Las limitaciones de la memoria a corto plazo}

La ``memoria a corto plazo'' del Agente LLM es la ventana de contexto (context window), que tiene tres limitaciones fundamentales:

\begin{enumerate}
    \item \textbf{Solo se puede anexar} (append-only): solo se pueden agregar tokens nuevos continuamente, no se pueden editar ni eliminar
    \item \textbf{Capacidad limitada}: incluso con una ventana de un millón de tokens, sigue habiendo un límite
    \item \textbf{No persiste}: al terminar la tarea o al vaciarse el contexto, toda la ``memoria'' desaparece
\end{enumerate}

\begin{warningbox}{El dilema del pez dorado}
Un Agente sin memoria a largo plazo es como un pez dorado: aunque hoy demuestre la conjetura de Riemann, mañana tiene que empezar desde cero. Peor aún, el siguiente intento no necesariamente tendrá éxito.
\end{warningbox}

\subsection{Reflexion: la forma más simple de memoria a largo plazo}

Reflexion es una extensión natural de ReAct:

\begin{enumerate}
    \item El Agente intenta resolver una tarea (por ejemplo, escribir código)
    \item La tarea falla (por ejemplo, no pasa las pruebas unitarias)
    \item El Agente \textbf{reflexiona} sobre la causa del fallo (por ejemplo, ``olvidé manejar el caso límite'')
    \item El resultado de la reflexión se escribe en la memoria a largo plazo
    \item En el siguiente intento se lee la memoria a largo plazo para evitar repetir el error
\end{enumerate}

\begin{knowledgebox}{Un nuevo paradigma de aprendizaje}
Reflexion es, en esencia, un método de \textbf{aprendizaje mediante lenguaje}, distinto del RL tradicional, que actualiza los pesos mediante descenso de gradiente:
\begin{itemize}
    \item No requiere una señal de recompensa escalar, puede aprovechar cualquier retroalimentación textual (errores de compilación, resultados de pruebas, comentarios de un instructor, etc.)
    \item No actualiza parámetros mediante retropropagación, sino que influye en el comportamiento futuro actualizando la memoria textual
\end{itemize}
\end{knowledgebox}

\subsection{Formas más complejas de memoria a largo plazo}

\begin{itemize}
    \item \textbf{Voyager}: aprende habilidades en Minecraft, guarda el código exitoso (por ejemplo, ``fabricar una espada'') en una biblioteca de habilidades que puede invocarse directamente en tareas posteriores.
    \item \textbf{Generative Agents}: 25 Agentes que simulan humanos viven en un pueblo; cada Agente mantiene un registro de eventos (episodic memory) y una memoria semántica (semantic memory) extraída de su experiencia, y estas memorias influyen en su comportamiento posterior.
    \item \textbf{Fine-tuning como memoria a largo plazo}: también se puede hacer fine-tuning directamente sobre los parámetros del modelo, ``escribiendo'' la experiencia en los pesos de la red neuronal.
\end{itemize}

\subsection{KOALA: una abstracción unificada del Agente}

Shunyu Yao (姚顺雨) propuso el marco KOALA, que describe de manera unificada cualquier Agente mediante tres componentes:

\begin{enumerate}
    \item \textbf{Memory}: el lugar donde se almacena la información (ventana de contexto + almacenamiento externo + parámetros del modelo)
    \item \textbf{Action Space}: el conjunto de acciones que el Agente puede ejecutar
    \item \textbf{Decision-making Procedure}: cómo elegir la siguiente acción dados la memoria y el espacio de acciones
\end{enumerate}

\subsection{Resumen de la sección}

La memoria a largo plazo hace que el Agente evolucione de ser un solucionador de un solo uso, ``sin estado'', a un sistema de aprendizaje capaz de acumular experiencia y mejorar de manera continua. La memoria puede ser un registro textual, una biblioteca de código de habilidades o incluso los propios parámetros del modelo.
